\section{转移轨道设计}
\label{sec:transfer}

e2m2e 的转移设计覆盖脉冲、低能量与低推力三类推进模式，核心方法论是
“搜索—优化”两步法：先在粗网格上以弹道传播筛选可行解族，再把最优候选
交给非线性规划（NLP）精化。统一的设计变量为切向速度比 $\alpha$
（控制出发方向）、转移时间 $T$ 与远地点到插入点的时间
$t_{\mathrm{ins}}$，总目标函数
\begin{equation}
  J = \Delta v_1 + \Delta v_2 ,
\end{equation}
约束包括插入点位置连续、速度方向与目标平行，以及对地球与月球的碰撞
避免。动力学基准为 CR3BP，质量比 $\mu$ 取自 DE421 常量基准。

\subsection{脉冲转移}

\subsubsection{Lambert 求解与 porkchop 扫描}

Lambert 求解器采用 Izzo (2015) 算法，内核在 Rust 侧（Python 仅做类型
转换与封装）：支持短路径/长路径方向与多圈解（多圈返回低能右分支），
并提供批量接口——$N$ 组几何 $\times\,M$ 个飞行时间的网格一次进入 Rust
求解，不可行组合填 NaN 而不影响其余。求解器经 Vallado 经典算例验证。

在出发时刻 $\times$ 飞行时间网格上逐点解 Lambert 即得 porkchop 双脉冲
$\Delta V$ 图，端点状态经终端条件抽象（\texttt{TerminalCondition}）提取，
解析终端无需动力学传播。

\subsubsection{三体打靶}

二体 Lambert 在地月空间只配当初猜。三体 Lambert 打靶用二体解播种 CR3BP
下的阻尼牛顿迭代：带状态转移矩阵 $\stm(t, t_0)$ 传播，每步修正量由末端
STM 的 $\stm_{rv}$ 块解出；全步使误差增大时步长减半；终端位置误差小于
$10^{-8}$（无量纲）判收敛。

\subsubsection{多脉冲优化与 Lawden 主矢量检验}

固定端点间的 $n$ 脉冲转移：决策变量仅为中途节点的 $[t_i, \vect{r}_i]$，
相邻节点间弧段由 Lambert 封闭（可切三体打靶精修），以 SLSQP 最小化总
$\Delta V$。双脉冲解是否还能改进，用 Lawden 主矢量检验回答：由横截条件
$\vect{p}(t_0) = \hat{\Delta v}_0$、$\vect{p}(t_f) = \hat{\Delta v}_f$
出发，协态经 STM 携载得到 $\vect{p}(t)$ 曲线，最优性必要条件为
\begin{equation}
  \norm{\vect{p}(t)} \le 1 \quad \text{全程成立},
  \qquad \text{脉冲点 } \norm{\vect{p}} = 1 \text{ 且与脉冲同向}.
\end{equation}
弧内 $\norm{\vect{p}} > 1$ 表明插入中途脉冲可进一步降低总 $\Delta V$
（Lion \& Handelsman, 1968），检验同时给出建议的插入时刻与位置，直接
播种三脉冲优化。LEO$\to$GEO 算例复现霍曼基准 $3.7708$ km/s。

此外，经典 Hohmann 双脉冲最小能量转移
\begin{equation}
  \Delta v_1 = \sqrt{\frac{\mu}{r_1}}\left(\sqrt{\frac{2 r_2}{r_1+r_2}} - 1\right),
  \qquad
  \Delta v_2 = \sqrt{\frac{\mu}{r_2}}\left(1 - \sqrt{\frac{2 r_1}{r_1+r_2}}\right)
\end{equation}
作为解析基准与复杂转移的初猜来源提供。

\subsection{低能量转移}

\subsubsection{月球借力（LGA）}

采用圆锥曲线拼接的三段式设计：双曲线脱离出发轨道 $\to$ 月球影响球内
借力改变速度矢量 $\to$ 双曲线插入目标轨道。搜索接口在 CR3BP 无量纲
状态下工作，典型搜索域为飞行时间 15--45 天、总 $\Delta v$ 上限
4.0 km/s，结果报告总代价、飞行时间与近月点高度。

\subsubsection{WSB 弹道捕获}

利用日地弱稳定边界（WSB）：航天器飞出月球影响球进入太阳摄动下的远端
弧段，远地点附近太阳引力做功，使月心相对速度低于捕获阈值，以近零
$\Delta v$ 弹道捕获（Belbruno \& Miller, 1993）。捕获判据为近月点处
二体 Kepler 能量
\begin{equation}
  H_2 = \frac{\norm{\vect{v}}^2}{2} - \frac{\mu_{\mathrm{moon}}}{r} < 0 .
\end{equation}
求解分两步：\textbf{弹道网格搜索}在 BCR4BP（双圆限制性四体问题）下前向
传播，于太阳相位 $\times$ 出发相位 $\times$ 飞行时间网格上按近月点高度
与 $H_2$ 判据筛选——传播、截面检测与筛选全部在 Rust 侧以 Rayon 并行
（Python 对照实现仅在显式指定时运行，绝不自动回退）；\textbf{到达段
精化}用三体 Lambert 打靶闭环最优候选的月心到达弧。典型搜索域为飞行
时间 90--150 天。

\subsubsection{不变流形拼接}

第三条低能量路线利用周期轨道自身的不变流形管：出发轨道的不稳定流形
与目标轨道的稳定流形在同一 Poincaré 截面（平面截面或近拱点截面
$s = \vect r \cdot \vect v = 0$）上的穿越点两两配对，按加权拼接代价
\begin{equation}
  w_r \norm{\Delta \vect r} + w_v \norm{\Delta \vect v}
\end{equation}
升序筛选候选。端到端流水线串联为：出发不稳定流形 $\times$ 目标稳定
流形（$\pm$ 分支四种组合全局取最优）传播到次天体近拱点截面取最优
拼接，拼接点之后由三体 Lambert 打靶闭合到目标轨道；脉冲相应分为出发、
拼接、到达三段。基准算例（L1 Lyapunov 族中等到大幅值轨道）的拼接脉冲
在几十 m/s 量级。当前该路线基于 CR3BP，星历接入已预留接口。

低能量路线的价值在于：以时间换燃料，为经济型着陆器/入轨器方案乃至应急
低燃料救援提供通道。

\subsection{低推力转移}

低推力段采用 7 维增广状态 $[\vect{r}, \vect{v}, m]$，动力学为
\begin{equation}
  \dot{\vect r} = \vect v, \qquad
  \dot{\vect v} = \vect a_{\mathrm{grav}} + \frac{T\,\delta}{m}\,\hat{\vect u}, \qquad
  \dot m = -\frac{T\,\delta}{I_{sp}\, g_0},
\end{equation}
其中油门 $\delta \in [0,1]$，推力方向 $\hat{\vect u}$ 由球面角
$(\theta_1, \theta_2)$ 参数化。求解是两级闭环：\textbf{Q-law 初猜}以
Lyapunov 反馈律前向积分产生次优控制历史；\textbf{SLSQP 打磨}在每段常值
控制 $(\delta, \theta_1, \theta_2)$ 上做最小燃料 NLP，终端位置速度匹配
目标。两种求解器可选：打靶法（解析雅可比，快 5--24 倍）与
Hermite--Simpson 配点法（大规模问题更鲁棒）。

脉冲与低推力的代价口径经等效 $\Delta v$ 区分：脉冲为
$\Delta v = \sum_i \norm{\Delta \vect v_i}$，低推力用齐奥尔科夫斯基反演
$\Delta v_{\mathrm{eq}} = I_{sp}\, g_0 \ln(m_0/m_f)$；由于引力损耗，
$\Delta v_{\mathrm{eq}} \ge$ 脉冲 $\Delta V$，大能量转移上差异显著。

\subsection{网格搜索 + NLP 两步优化框架}

第一步\textbf{网格搜索}：在出发轨道上等时间采样出发点（典型 200 个），
每个出发点在 $\alpha \in [0.5, 2.5]$ 内均匀采样（典型 101 个），前向
积分转移弧，经碰撞（地球 200 km / 月球 100 km）、相交与最小距离
（$10^{-3}$）约束筛选可行解；数值内核已下沉 Rust 并行执行。

第二步 \textbf{NLP 精化}：决策变量 $\vect{y} = \{\alpha, T, t_{\mathrm{ins}}\}$，
目标 $J(\vect y) = \Delta v_1 + \Delta v_2$，约束为插入点位置连续
（等式）与速度平行（等式，或松弛为夹角容差 0.05 rad 的不等式）。求解器
默认 SciPy SLSQP，可选 COPT 商业优化器并在失败时回退。

\begin{table}[htbp]
\centering\small
\caption{转移设计能力一览}
\label{tab:transfer}
\begin{tabularx}{\textwidth}{@{}llX@{}}
\toprule
\textbf{类别} & \textbf{方法} & \textbf{要点} \\
\midrule
脉冲 & Lambert（Izzo） & Rust 内核，批量网格求解，短/长路径与多圈 \\
脉冲 & 三体打靶 & 阻尼牛顿 + STM，位置误差 $< 10^{-8}$ 收敛 \\
脉冲 & 多脉冲 & SLSQP + Lawden 主矢量检验，自动建议插入脉冲 \\
脉冲 & Hohmann & 解析基准与初猜 \\
低能量 & LGA & 圆锥曲线拼接三段式，TOF 15--45 天 \\
低能量 & WSB & BCR4BP + $H_2<0$ 判据，Rust Rayon 并行搜索 \\
低能量 & 流形拼接 & Poincaré 截面配对 + 三体打靶闭合，拼接脉冲 m/s 级 \\
低推力 & Q-law + NLP & 打靶（解析雅可比 5--24$\times$）或配点 \\
框架 & 搜索--优化 & 网格搜索（Rust 并行）$\to$ SLSQP/COPT 精化 \\
\bottomrule
\end{tabularx}
\end{table}
